% !TEX root = ../../tree_notes.tex
% =====================================================================
%  9.1  The Cartesian Tree in O(n)              sec:c26-cartesian
%  Exports: def:c26-cartesian  thm:c26-unique  lem:c26-span
%           fig:c26-build  lst:c26-stackbuild  ex:c26-cart-run
%           prob:c26-goodsub  lst:c26-goodsub-soln   (practice: LC 907)
% =====================================================================
\section{The Cartesian Tree in $O(n)$}\label{sec:c26-cartesian}

\begin{previouslybox}
\begin{enumerate}[nosep, leftmargin=1.6em]
  \item A binary tree's \emph{inorder} walk visits \text{left}, node, \text{right}. On a BST it lists
        the keys sorted. Which property of the tree makes that true? (\chref{c03}, \thmref{c05-inorder})
  \item A \emph{min-heap} stores its smallest element where, and what is the parent--child rule?
        (at the root; every parent $\le$ its children -- the idea Chapter~10 builds, \chref{c07})
  \item You have a \code{stack}. Pushing and popping a sequence once costs what in total, and why can
        an inner \code{while}-pop loop still be $\bigO{n}$ overall? (each element pushed/popped once)
  \item Given an array, its minimum sits at some position $p$. What two pieces does $p$ split the
        array into? (the strictly-left part $a[1..p-1]$ and strictly-right part $a[p+1..n]$)
\end{enumerate}
\end{previouslybox}

\begin{topicintro}
There is a tree sleeping inside every array. To wake it, impose \emph{two} ordering rules at once: the
tree's inorder reading must be the array left-to-right (so the \emph{index} behaves like a BST key),
and the tree must be a heap on the array's \emph{values} (so a parent beats its children). A structure
that is simultaneously a BST on position and a heap on value is the \term{Cartesian tree}. On distinct
values it is unique (\thmref{c26-unique}), and its root is nothing but the position of the minimum.

Two facts make it worth your attention. First, it is \emph{buildable in }$\bigO{n}$ by a single
monotonic-stack sweep -- you never touch the $O(n\log n)$ recursive-minimum route. Second, and this is
the chapter's beating heart: the subtree hanging under position $i$ is \emph{exactly} the maximal
window over which $a[i]$ is the minimum (\lemref{c26-span}). Every "for each element, how many ranges
does it win?" problem is therefore already solved once you can read a subtree, and every such problem
becomes two stack passes. We build the tree, prove it unique, state the span lemma, and cash it in on
LeetCode 1793 -- find the best subarray through a fixed index $k$.
\textbf{Prerequisites.} Arrays and the $\bigO{n}$ amortised stack discipline; the inorder walk from
Chapter~3 (\chref{c03}) as a black box.

\textbf{After this section you can}
\begin{itemize}[nosep]
  \item write down the two invariants defining a Cartesian tree and build one on a small array by hand; \hfill\lv{3}
  \item prove the tree is unique when values are distinct, and explain what breaks with ties; \hfill\lv{3}--\lv{4}
  \item build it in $\bigO{n}$ with a monotonic stack and read off every node's parent; \hfill\lv{4}
  \item use the span-equals-subtree lemma to turn a "max over ranges containing $k$" into per-index spans. \hfill\lv{4}--\lv{5}
\end{itemize}
\end{topicintro}

\begin{recognitionbox}[Pattern card -- the array \emph{is} a tree]
\reckey{a question about \emph{each element's dominance window} -- where it is the min/max before a
smaller/larger element appears -- or a \emph{fixed index that every candidate range must contain}.}
\recuse{
  \item "largest rectangle through each bar", "sum of subarray minimums", "how many ranges have $a[i]$
        as their smallest" -- per-index spans, \lemref{c26-span};
  \item "subarray with best $\min \cdot \text{length}$ that must contain $k$" -- \probref{c26-goodsub};
  \item "rebuild a max-tree from an array" -- the same object with the heap rule flipped to max,
        \secref{c26-treapview}.
}
\recnotuse{
  \item the array is \emph{changing} between queries -- you need a segment tree (\chref{c13}), not a
        tree read off a frozen snapshot;
  \item you need the \emph{global} min/max only -- one variable, no structure;
  \item you want a \emph{balanced} tree you can insert into for $\bigO{\log n}$ -- that is Chapter~8's
        treap (\chref{c06}); here the shape is dictated by values, not by you.
}
One-line test: \emph{can the array be read as ``a BST by index, a heap by value'' at once? Then its
Cartesian tree, and the span lemma, are waiting to be used.}
\end{recognitionbox}

\subsection{Intuition}

Take the running array $A = (3,1,4,2)$ from the opener. Ask: \emph{which element should be the root of
a tree that is a heap on values?} A heap forces the smallest element to the top, so the root must be
$1$, sitting at position $2$. Now the BST-by-index rule says everything with index $<2$ goes to the
left of the root and everything with index $>2$ goes to its right -- so $a_1=3$ becomes the entire left
subtree, and the pair $a_3=4,a_4=2$ becomes the right subtree. Inside that right subtree the smallest is
$2$ (position $4$), so it is the local root and $4$ (position $3$, index $<4$) lands on its \emph{left}.
That is the whole tree in \figref{c26-cartA}, and it is the same figure \figref{c26-build} builds step
by step: \emph{root $=$ the minimum, then recurse on the two halves the minimum splits off}.

The remarkable coincidence -- and the seed of every application in this chapter -- is what the
subtrees \emph{mean}. Node $2$ (position $4$) owns a subtree covering just positions $3$--$4$, which is
precisely the widest window $(2,5)$ around position $4$ in which the value $2$ is the minimum: to its
left, position $2$ holds $1<2$ and blocks it; to its right there is nothing, so it runs to the end. The
subtree is the dominance window. Read that backwards and building the tree is the same act as finding,
for every position, its nearest smaller neighbours -- which is what a monotonic stack does in one pass.

\begin{figure}[H]
\centering
\begin{tikzpicture}[diagram, scale=0.92, transform shape]
  % the min-Cartesian tree of A = 3,1,4,2 with parent links highlighted
  \node[troot] (n2) at (0,0)  {$1$};
  \node[tleaf] (n1) at (-2.2,-1.5) {$3$};
  \node[tnode] (n4) at (2.2,-1.5) {$2$};
  \node[tleaf] (n3) at (1.1,-3.0) {$4$};
  \draw[tedgehl] (n2) -- (n1) node[midway, above, font=\scriptsize, text=darkgray] {$a_1$};
  \draw[tedgehl] (n2) -- (n4) node[midway, above, font=\scriptsize, text=darkgray] {$a_4$};
  \draw[tedgehl] (n4) -- (n3) node[midway, right, font=\scriptsize, text=darkgray] {$a_3$};
  \node[annot, align=left, font=\scriptsize] at (0,1.3) {root $=$ minimum of $A$, at position $2$};
  % array strip
  \node[font=\scriptsize] at (-3.0,-3.9) {$A=$};
  \foreach \x/\v in {-2.4/3,-1.6/1,-0.8/4,0.0/2}
    \node[cellnode] at (\x,-3.9) {\v};
  \node[annot, align=left, font=\scriptsize] at (3.4,-3.0)
    {\textbf{inorder} gives $3,1,4,2$\\ \textbf{heap} gives $1\le3,\,1\le2\le4$};
\end{tikzpicture}
\caption{The min-Cartesian tree of $A=(3,1,4,2)$, built by "root is the minimum, recurse on the two
halves". Position subscripts are shown on the edges. The inorder traversal reads the array back
($3,1,4,2$); every parent is $\le$ its children. This is \defref{c26-cartesian}; \thmref{c26-unique}
says it is the \emph{only} such tree, and \lemref{c26-span} says each subtree is that node's
dominance window.}\label{fig:c26-build}
\end{figure}

\subsection{Definitions and Notation}

\begin{definition}{Cartesian tree (min-heap form)}{c26-cartesian}
A \term{Cartesian tree} of an array $a[1..n]$ is a binary tree $T$ on the index set $\{1,\dots,n\}$
satisfying two conditions simultaneously:
\begin{itemize}[nosep]
  \item \textbf{BST order on the index.} The \emph{inorder} traversal of $T$ visits the indices in
        increasing order $1,2,\dots,n$ -- equivalently, the indices in any subtree form a contiguous
        interval, and a node's left subtree holds smaller indices, its right subtree larger ones.
  \item \textbf{Heap order on the value.} For every node $i$ and every child $j$ of $i$,
        $a[i] \le a[j]$. (A \emph{max-}Cartesian tree flips this to $a[i]\ge a[j]$.)
\end{itemize}
Equivalently, defined recursively: if $p$ is a position of the minimum of $a$, then $T$ has root $p$,
left subtree $=$ Cartesian tree of $a[1..p-1]$, right subtree $=$ Cartesian tree of $a[p+1..n]$.
\end{definition}

\noindent\textbf{In plain words.} Put the smallest value on top; everything to its left in the array
becomes the left branch, everything to its right the right branch; repeat inside each side. The array
order is the tree's inorder; the value order is the heap.

\begin{theorem}{Uniqueness under distinct values}{c26-unique}
If the $n$ values of $a$ are pairwise distinct, its min-Cartesian tree is unique. If values repeat,
the tree need not be unique: two equal minima at different positions can be ordered either way, and a
tie-break convention (which equal element becomes the ancestor) must be fixed to restore uniqueness.
\end{theorem}
\begin{proofbox}
By induction on $n$. The heap rule forces the root to be a position $p$ with the minimum value, because
the root must be $\le$ every node in its subtree, which is the whole array. With distinct values that
$p$ is unique. The BST (inorder) rule then forces all indices $<p$ into the left subtree (a contiguous
block) and all indices $>p$ into the right, since inorder lists indices ascending and $p$ sits in the
middle. The left and right blocks are independent arrays, so by the induction hypothesis each has a
unique Cartesian tree, and gluing them under root $p$ is forced. Hence the whole tree is unique. With
two equal minima $p,q$ either may serve as root of that combined subtree, giving two valid trees -- the
tie must be broken externally.
\end{proofbox}

\begin{lemma}{Subtree $=$ dominance window}{c26-span}
Let $L_i$ be the nearest index strictly left of $i$ with $a[L_i] < a[i]$ (or $0$ if none), and $R_i$ the
nearest index strictly right of $i$ with $a[R_i] < a[i]$ (or $n{+}1$ if none). Then the subtree of node
$i$ in the min-Cartesian tree spans exactly the open window $(L_i, R_i)$: the set
$\{L_i+1,\dots,R_i-1\}$. Consequently $a[i]$ is the minimum of that window, and the number of
subarrays for which $a[i]$ is \emph{the} minimum element read at position $i$ is
$(i-L_i)\cdot(R_i-i)$.
\end{lemma}
\begin{proofbox}
An ancestor of $i$ must have value $\le a[i]$ (heap rule) and lie outside $i$'s subtree, and the BST
rule makes the nearest such blockers exactly the closest \emph{smaller} values on each side, $L_i$ and
$R_i$: their positions are the highest ancestors of $i$ (they became $i$'s parent-chain because they
are the first values that could sit above it). No index inside $(L_i,R_i)$ has value $<a[i]$ by
definition of $L_i,R_i$, so $i$ is the minimum of that whole window and its subtree is exactly the
window (an ancestor of $i$ has a smaller value and lies at $L_i$ or $R_i$, outside it). The window
$(L_i,R_i)$ contains $i$; a subarray inside it that includes $i$ starts anywhere in $[L_i+1..i]$ (that
many $= i-L_i$ choices) and ends anywhere in $[i..R_i-1]$ ($R_i-i$ choices), giving the product. \end{proofbox}

\begin{insightbox}
\lemref{c26-span} is the hinge of the entire chapter. It says the geometric fact ``size of $i$'s
subtree'' and the combinatorial fact ``number of ranges where $a[i]$ is the champion'' are the
\emph{same thing}. So any sum of the form $\sum_i a[i]\cdot(\text{ranges where } a[i] \text{ wins})$ --
subarray minimums, the histogram, the Summit -- needs only the two arrays $L,R$, and those fall out of
one monotonic stack pass per side. Read the stack pass right, and you have built the tree; read it as
counts, and you have solved the problem. \secref{c26-monostack} makes the stack pass explicit; \secref{c26-rmq} reads the tree as an index.
\end{insightbox}

\subsection{The algorithm: one stack, one pass}

The recursive definition is $\bigO{n^2}$ in the worst case (a sorted array picks the minimum at one end
every time). The linear build instead walks the array left to right maintaining the \emph{right spine}
of the partially built tree -- the chain from the root down through right children -- in a stack. Values
on that spine are increasing; a new element $a[i]$ pops the spine until it finds its parent, and the
last node it popped becomes its left child.

\begin{lstlisting}[caption={The $O(n)$ monotonic-stack build of a min-Cartesian tree; \code{Lc[], Rc[]} are the left and right children}, label={lst:c26-stackbuild}]
// indices 1..n; a[1..n] the array. Builds children of the min-Cartesian tree.
vector<int> Lc(n+1, 0), Rc(n+1, 0), par(n+1, 0);
vector<int> st;                       // right spine: indices with increasing values
for (int i = 1; i <= n; ++i) {
    int last = 0;                     // last subtree popped -- becomes i's LEFT child
    while (!st.empty() && a[st.back()] > a[i]) {
        last = st.back(); st.pop_back();
    }
    if (!st.empty()) {                // current top beats a[i]: i hangs off it
        Rc[st.back()] = i;            //   as that node's right child
        par[i] = st.back();
    }
    if (last) {                       // the popped chain becomes i's left child
        Lc[i] = last;
        par[last] = i;
    }
    st.push_back(i);
}
int root = st.front();                // the bottom of the stack never got a parent: the root
\end{lstlisting}
\begin{complexitybox}
Time $\bigO{n}$: each index is pushed exactly once and popped at most once across the whole sweep, so
the inner \code{while} performs at most $n$ pops in total -- the classic amortised-stack argument.
Memory $\bigO{n}$ for the three child/parent arrays and the stack. Verdict: the $\bigO{n}$ build, versus
the $\bigO{n^2}$ "scan for the minimum" recursion (a linked-list array is its worst case). The build is
also, in disguise, the histogram stack of \secref{c26-monostack} -- here the stack \emph{is} the tree.
\end{complexitybox}

\subsection{Worked Example}

\begin{example}{Building the tree on $A=(3,1,4,2)$}{c26-cart-run}
Process left to right, tracking the stack $\mathit{st}$ and the children set so far.\\
\emph{$i=1$ ($a=3$).} Stack empty; nothing popped; no parent. Push $1$. $\mathit{st}=[1]$.\\
\emph{$i=2$ ($a=1$).} $a_1=3>1$: pop it, $\mathit{last}=1$. Stack empty, so no right-child link.
Set $\mathit{Lc}[2]=1$ (and $\mathit{par}[1]=2$). Push $2$. $\mathit{st}=[2]$.\\
\emph{$i=3$ ($a=4$).} $a_2=1>4$? no. Nothing popped. $\mathit{Rc}[2]=3$, $\mathit{par}[3]=2$. Push
$3$. $\mathit{st}=[2,3]$.\\
\emph{$i=4$ ($a=2$).} $a_3=4>2$: pop $3$, $\mathit{last}=3$. Now $a_2=1>2$? no. $\mathit{Rc}[2]=4$
(this \emph{overwrites} the earlier link, moving node $3$'s parent to $4$), $\mathit{par}[4]=2$; then
$\mathit{Lc}[4]=3$, $\mathit{par}[3]=4$. Push $4$. $\mathit{st}=[2,4]$.\\
Root $=\mathit{st.front()}=2$. Children: $\mathit{Lc}[2]=1,\ \mathit{Rc}[2]=4,\ \mathit{Lc}[4]=3$ --
exactly \figref{c26-build}: root $1$, left $3$, right $2$ with left child $4$. Inorder
$3,1,4,2$; min-heap throughout. \checkmark
\end{example}
\begin{tracebox}[Span of every node from \lemref{c26-span} on $A$]
Prev/next strictly-smaller $(L_i,R_i)$ with $A=(3,1,4,2)$: node $1$ ($a=3$): $L=0,R=2$; node $2$
($a=1$): $L=0,R=5$ (nothing smaller); node $3$ ($a=4$): $L=4$? no -- nearest smaller left is $a_2=1$
at $2$, nearest smaller right is $a_4=2$ at $4$, so $L=2,R=4$; node $4$ ($a=2$): $L=2,R=5$. Read off
the subtree spans: node $2$'s window $(0,5)$ is all four indices (its subtree is everything); node
$4$'s window $(2,5)$ is $\{3,4\}$ -- matching that $a_4{=}2$ is the minimum of exactly $a[3..4]$.
\end{tracebox}

\subsection{Variations and Twists}

\begin{itemize}[nosep]
  \item \textbf{Max-Cartesian tree}: flip the pop test to $a[\mathit{st.back()}] < a[i]$ and the root
        is the maximum -- this is exactly LeetCode's ``maximum binary tree'' (\secref{c26-treapview}). \hfill\lv{3}
  \item \textbf{Building by the divide route}: recursively find the min (via an RMQ index from
        \secref{c26-rmq}) -- $\bigO{n}$ if you already have a linear build to seed the RMQ, else
        $\bigO{n^2}$. The stack build dominates it. \hfill\lv{4}
  \item \textbf{Parent from one neighbour}: the parent of $i$ is whichever of $L_i, R_i$ has the
        \emph{larger} value (the later of the two blockers to have become a spine node) -- handy for
        reconstructing without the stack. \hfill\lv{4}
  \item \textbf{Ties}: with duplicates, decide whether $L/R$ use $<$ or $\le$ on each side; the same
        asymmetry that prevents double counting in the Summit (\probref{c26-summit}). \hfill\lv{5}
\end{itemize}

\subsection{Common Pitfalls}\label{sec:c26-cartesian-pitfalls}

\begin{warningbox}[Trap 1 -- trusting the recursive minimum scan]
Implementing the definition literally (scan the interval for the minimum, recurse) is $\bigO{n^2}$ on
an already-sorted array -- the minimum always sits at one end, the recursion is a chain of $n$ scans.
For $n=10^5$ that is $\sim 10^{10}$ and the build itself \tle s. Use the stack build
\coderef{c26-stackbuild}; its amortised argument is exactly the "push once, pop at most once" fact.
\end{warningbox}
\begin{lstlisting}[caption={\textcolor{warningred}{WRONG}: the naive divide-by-minimum build is $O(n^2)$ on sorted input}, language=C++, basicstyle=\ttfamily\small]
Node* naive(int lo, int hi) {          // worst case: array sorted -> depth n, n per level
    if (lo > hi) return nullptr;
    int p = argmin(a, lo, hi);         // O(hi-lo+1) SCAN at every node  => O(n^2) total
    Node* r = new Node(p);
    r->l = naive(lo, p-1); r->r = naive(p+1, hi);
    return r;
}
\end{lstlisting}
\begin{lstlisting}[caption={\textcolor{successgreen}{RIGHT}: the monotonic stack builds it once, in $O(n)$}, language=C++, basicstyle=\ttfamily\small]
// the sweep of lst:c26-stackbuild: one push and at most one pop per index.
// No scanning, no recursion, no degenerate-depth blow-up.
for (int i=1;i<=n;i++){ while(!st.empty() && a[st.back()]>a[i]){last=st.back();st.pop_back();}
    if(!st.empty()) Rc[st.back()]=i; if(last) Lc[i]=last; st.push_back(i); }
\end{lstlisting}

\begin{warningbox}[Trap 2 -- forgetting the tie rule]
\thmref{c26-unique} guarantees uniqueness \emph{only for distinct values}. If $a$ repeats, popping on
$>$ versus $\ge$ silently picks a different tree and changes which window each element ``owns''. For a
\emph{max} problem (like \probref{c26-goodsub}) any consistent rule works because you take a maximum
over all positions; for a \emph{sum} problem (LC 907, the Summit) you must break ties asymmetrically --
strict on one side, non-strict on the other -- or a range's minimum gets counted twice.
\end{warningbox}

\begin{warningbox}[Trap 3 -- mixing up which order is which]
The BST order is on the \emph{index} and the heap order is on the \emph{value}. Swapping them gives the
``insertion-order tree'' of \secref{c26-treapview}, a different (also useful) object. The inorder that
reconstructs the array uses the index order -- if your inorder prints something other than
$a[1],\dots,a[n]$ you have a heap-shaped BST, not a Cartesian tree.
\end{warningbox}

\subsection{Problem Lab}

The section's anchor is LeetCode 1793 -- among all subarrays that \emph{contain a fixed index} $k$,
maximise $(\text{minimum of the subarray}) \times (\text{its length})$. The naive double loop is
$\bigO{n^2}$; the span lemma collapses it to two stack passes and one scan, because the only subarrays
worth checking are the full dominance windows of \lemref{c26-span}.

% ---------------------------------------------------------------------------
\solvedproblem{c26-goodsub}{Maximum Score of a Good Subarray}{LeetCode 1793}
  {https://leetcode.com/problems/maximum-score-of-a-good-subarray/}{LC Hard · L5}
  {Cartesian span lemma, prev/next smaller via monotonic stack, maximise min times width}

\begin{problemstatement}[Maximum Score of a Good Subarray]
You are given a $0$-indexed integer array $\mathit{nums}$ and an integer $k$. A \textbf{good} subarray
$(i, j)$ is one with $i \le k \le j$. Its \textbf{score} is
$\min(\mathit{nums}[i], \dots, \mathit{nums}[j]) \cdot (j - i + 1)$. Return the \textbf{maximum} score
among all good subarrays. $1 \le \mathit{nums.\text{length}} \le 10^{5}$, $1 \le \mathit{nums}[i] \le 10^{5}$,
and $0 \le k < \mathit{nums.\text{length}}$.
\end{problemstatement}

\begin{attemptgate}[Attempt it yourself first -- \lv{5}, give yourself 30 minutes]
The double loop over all $(i,j)$ through $k$ is $\bigO{n^2}$. The winning move is to stop thinking about
subarrays and think about \emph{which element is the minimum of the chosen subarray}.
\peek{Fix a position $m$. What is the \emph{longest} subarray through $k$ on which $\mathit{nums}[m]$
is the minimum? (Hint: the dominance window of \lemref{c26-span}.)}
\nudge{For $\mathit{nums}[m]$ to be the minimum, the subarray cannot cross a strictly-smaller element,
so it lies inside $(L_m, R_m)$. If that window contains $k$, the best subarray with minimum at $m$ is the
whole window, scoring $\mathit{nums}[m]\cdot(R_m - L_m - 1)$. Compute every $L_m, R_m$ with two stack passes.}
\fullhint{Take $\max$ over $m$ whose window contains $k$ of $\mathit{nums}[m]\cdot(R_m-L_m-1)$, where
$L_m$ ($R_m$) is the nearest index to the left (right) with a strictly smaller value. Any good subarray
has its minimum at some $m$, and its full window through $m$ is at least as good and still contains $k$.
The answer for the sample is $3 \times 5 = 15$.}
\end{attemptgate}

\solstep{Understanding the Problem}
We maximise $(\min \text{ over a window}) \times (\text{window length})$ subject to the window
containing index $k$. The two factors pull opposite ways -- a longer window is smaller in minimum -- so
naively you must try every window through $k$, $\bigO{n^2}$. The reframing: the value $\min$ is attained
at \emph{some} position $m$ in the window, so enumerate $m$ instead of the window. For a fixed $m$,
the minimum over any window containing $m$ is at most $\mathit{nums}[m]$, and equals it as long as the
window stays inside $(L_m, R_m)$ -- the dominance window of \lemref{c26-span}, found by nearest
strictly-smaller neighbours. So the \emph{best window that takes $\mathit{nums}[m]$ as its minimum} is
the whole $(L_m, R_m)$, \emph{if} that window reaches across $k$.

\solstep{Key Observations}
\begin{observation}
For a good subarray with score $s$, let $m$ be a position attaining its minimum. The dominance window
$(L_m, R_m)$ contains that subarray (nothing inside it is strictly smaller than $\mathit{nums}[m]$),
hence contains $k$ too, and using the whole window keeps the minimum $\mathit{nums}[m]$ while never
shrinking the length. So it is never worse to expand to the full window.
\end{observation}
\begin{observation}
Therefore the optimum equals $\max$ over $m$ with $L_m < k < R_m$ of
$\mathit{nums}[m]\cdot(R_m - L_m - 1)$. Only windows that \emph{straddle} $k$ are candidates.
\end{observation}
\begin{observation}
Both $L_m$ and $R_m$ (nearest strictly-smaller on each side) come from monotonic-stack passes -- one
left-to-right, one right-to-left -- in $\bigO{n}$. That is the whole algorithm.
\end{observation}

\solstep{Think It Through With Me}
\begin{thinkbox}
\deadend{Enumerate $i \le k$ and $j \ge k$ and compute the running minimum.}
{$\bigO{n^2}$: for each start you grow the end and update the min -- on $n=10^5$ that is $\sim 10^{10}$,
a \tle. It also hides the structure; every extra $j$ beyond the dominance window contributes nothing new.}
\deadend{Binary-search the answer on the minimum value.}
{Works ($\bigO{n\log V}$: is there a length-$k$-containing window with all values $\ge$ a threshold?),
but it is a heavier hammer than needed and obscures \emph{why} the answer is a per-position span.}
\aha{Assign each good subarray's score to the \emph{position of its minimum}. That position's dominance
window is the best subarray for that minimum; keep only windows that cover $k$. Two stack passes give
every window, one scan gives the max -- $\bigO{n}$ and a direct hit on \lemref{c26-span}.}
\end{thinkbox}

\solstep{Solution}
Compute $L[m]$ = nearest index $< m$ with $\mathit{nums}[L]<\mathit{nums}[m]$ (else $-1$) with an
increasing stack scanning left-to-right; compute $R[m]$ = nearest index $> m$ with
$\mathit{nums}[R]<\mathit{nums}[m]$ (else $n$) by scanning right-to-left the same way. Then the window
in which $\mathit{nums}[m]$ is the minimum is $(L[m], R[m])$, length $R[m]-L[m]-1$. It is a good
subarray iff $L[m] < k < R[m]$. Answer is the maximum of $\mathit{nums}[m]\cdot(R[m]-L[m]-1)$ over such
$m$. Because we only ever take a maximum over all positions, a consistent choice of $<$ on both sides is
safe -- duplicates cannot cause an under-count here.

\solstep{Proof of Correctness}
\begin{proofbox}[The span max equals the true optimum]
(\emph{Achievable}) For each $m$ with $L[m]<k<R[m]$, the subarray $[L[m]+1,\, R[m]-1]$ contains $k$
(hence is good) and has minimum exactly $\mathit{nums}[m]$ (no strictly smaller value lies inside the
open window, and $\mathit{nums}[m]$ itself is in it). So its score $\mathit{nums}[m]\cdot(R[m]-L[m]-1)$
is attainable and the reported maximum is a valid lower bound on the optimum. (\emph{Optimal}) Let
$[i,j]$ be any good subarray of maximal score and $m$ a position of its minimum. Since all values in
$[i,j]$ are $\ge \mathit{nums}[m]$, no index of $[i,j]$ has a value strictly below $\mathit{nums}[m]$,
so $[i,j]\subseteq (L[m], R[m])$; as $[i,j]\ni k$ we get $L[m]<k<R[m]$. The window
$[L[m]+1,R[m]-1]$ contains $[i,j]$, so it is at least as long with the same minimum, giving score
$\ge$ that of $[i,j]$. Thus some candidate considered by the algorithm is $\ge$ the optimum's score --
combined with achievability, the maximum the algorithm returns \emph{is} the optimum. \end{proofbox}

\solstep{Implementation}
\begin{lstlisting}[caption={Maximum Score of a Good Subarray -- spans from two stack passes (LeetCode 1793)}, label={lst:c26-goodsub-soln}]
class Solution {
public:
    int maximumScore(vector<int>& nums, int k) {
        int n = (int)nums.size();
        vector<int> L(n, -1), R(n, n), st;
        for (int i = 0; i < n; ++i) {                 // prev strictly smaller
            while (!st.empty() && nums[st.back()] >= nums[i]) st.pop_back();
            if (!st.empty()) L[i] = st.back();
            st.push_back(i);
        }
        st.clear();
        for (int i = n - 1; i >= 0; --i) {            // next strictly smaller
            while (!st.empty() && nums[st.back()] >= nums[i]) st.pop_back();
            if (!st.empty()) R[i] = st.back();
            st.push_back(i);
        }
        long long ans = 0;
        for (int m = 0; m < n; ++m)
            if (L[m] < k && k < R[m])                 // its window straddles k
                ans = max(ans, 1LL * nums[m] * (R[m] - L[m] - 1));
        return (int)ans;
    }
};
\end{lstlisting}

\solstep{Trace}
\begin{tracebox}[Official sample: \texttt{nums = [1,4,3,7,4,5]}, \texttt{k = 3}]
Prev-smaller $L$: $[-1,0,0,2,2,4]$. Next-smaller $R$: $[6,2,6,4,6,6]$. Keep $m$ with $L[m]<3<R[m]$:
$m=0$: $1\cdot(6-(-1)-1=6)=6$. $m=2$: $L=0<3<R=6$, $3\cdot5=15$. $m=3$: $L=2<3<R=4$, $7\cdot1=7$.
$m=4$: $L=2<3<R=6$, $4\cdot3=12$. $m=1,5$ fail ($R[1]=2\not>3$, $L[5]=4\not<3$). Maximum $=15$ --
the subarray $[1..5]=[4,3,7,4,5]$, minimum $3$, length $5$, score $15$. \checkmark matches the official
output.
\end{tracebox}

\solstep{Complexity}
\begin{complexitybox}
Time $\bigO{n}$: two monotonic-stack passes (each index pushed once, popped at most once) plus one
$\bigO{n}$ scan for the maximum. Memory $\bigO{n}$ for $L,R$ and the stack. Against the $\bigO{n^2}$
window enumeration this is the full win the span lemma promises -- and it is literally the Cartesian
build \coderef{c26-stackbuild} run once per side. Verdict: \ac.
\end{complexitybox}

\solstep{Pitfalls and Variants}
\begin{warningbox}
Three ways to fall: (i) treating a position whose window does \emph{not} contain $k$ as a candidate --
it is not a good subarray, skip it (the $L[m]<k<R[m]$ guard); (ii) using $\le$ instead of $<$ so a
duplicate minimum truncates the window early -- harmless for a pure \emph{max}, but know that you chose
it; (iii) $32$-bit overflow in $\mathit{nums}[m]\cdot\text{length}$ ($10^5 \cdot 10^5 = 10^{10}$) --
multiply in \code{long long}. \emph{Variant:} the two-pointer "grow toward the taller neighbour" greedy
also passes, but the span proof is the one that generalises to the Summit's signed sum.
\end{warningbox}

\begin{insightbox}[Key takeaways]
\begin{enumerate}[nosep, leftmargin=1.6em]
  \item \textbf{A Cartesian tree is a BST on index and a heap on value at once} (\defref{c26-cartesian}),
        unique on distinct values, root $=$ the minimum (\thmref{c26-unique}).
  \item \textbf{The subtree is the dominance window} (\lemref{c26-span}) -- the pivot from a tree-shape
        fact to a counting fact, and the reason one stack pass solves so much.
  \item \textbf{Build it, or just read it.} \coderef{c26-stackbuild} materialises the tree; for scoring
        problems you keep only $L,R$ and never link a single node (\probref{c26-goodsub}).
\end{enumerate}
\end{insightbox}

\begin{hardeningbox}
\begin{itemize}[nosep, leftmargin=1.4em]
  \item \textbf{Best subarray through each index} -- run the same span max for every $k$; the $L,R$
        arrays answer all $k$ after one build.
  \item \textbf{Largest all-positive rectangle under a skyline} -- the histogram of
        \secref{c26-monostack} is \emph{exactly} this span maximum with width forced to start $=1$; the
        link is made explicit there.
  \item \textbf{Count the ranges, don't max them} -- replace $\max$ by a sum and you are on the road to
        LC 907 and the Summit; the only new care is asymmetric tie-breaking.
\end{itemize}
\end{hardeningbox}

\subsection{Practice Problems}

\begin{stuckbox}
\begin{itemize}[nosep, leftmargin=1.4em]
  \item Did you assign each subarray's value to the \emph{position of its minimum}, not to the window?
        (\probref{c26-goodsub}.)
  \item Did you guard that the window must \emph{contain} $k$ before scoring it?
  \item For the sum version: did you make one side strict and the other non-strict so a range with two
        equal minima is counted once, not twice?
\end{itemize}
\end{stuckbox}

\begin{practiceset}
  \practice{Sum of Subarray Minimums}{LeetCode 907}
    {https://leetcode.com/problems/sum-of-subarray-minimums/}{LC Medium · L3}
    {Return $\sum$ over all subarrays of their minimum, modulo $10^9+7$; $\mathit{arr}.\text{length}\le 4\cdot10^4$.
    This is \lemref{c26-span} turned from a maximum into a sum: the answer is
    $\sum_i a[i]\cdot(i-L_i)\cdot(R_i-i)$, the count of subarrays whose minimum is $a[i]$ read at $i$,
    from one stack pass for each of $L_i, R_i$. The one trap \emph{not} present in \probref{c26-goodsub}:
    a subarray with two equal minima must be credited to exactly one position, so break ties by taking
    \emph{strict} $<$ on one side (say prev smaller) and \emph{non-strict} $\le$ on the other (next
    smaller-or-equal). Keep every partial product in \code{long long} and reduce modulo $10^9+7$. The
    sample $[3,1,2,4]$ gives $17$ -- and the same two-pass-with-asymmetric-ties template is exactly what
    the Summit (\probref{c26-summit}) will run for both the max and the min side at once.}
\end{practiceset}

\begin{recallbox}[10-second refresher -- 9.1 The Cartesian Tree in $O(n)$]
A Cartesian tree is \emph{inorder $=$ index} and \emph{heap $=$ value}: root is the minimum, unique on
distinct values (\defref{c26-cartesian}, \thmref{c26-unique}). One monotonic stack builds it in
$\bigO{n}$ by maintaining the increasing right spine (\coderef{c26-stackbuild}). Its key payoff,
\lemref{c26-span}, says node $i$'s subtree is the maximal window where $a[i]$ is the minimum -- so
"each element's dominance range" problems need only the prev/next smaller arrays. LeetCode 1793
(\probref{c26-goodsub}) maximises $\min\times\text{length}$ through $k$ by scoring exactly those windows;
LC 907 sums them once you break ties asymmetrically.
\selftest{
  \item Why must the root of a min-Cartesian tree be the position of the minimum?
  \item State \lemref{c26-span} in one sentence and read $a_4$'s window off $A=(3,1,4,2)$.
  \item Why is the naive divide-by-minimum build $\bigO{n^2}$, and what input forces it?
  \item In \probref{c26-goodsub}, why is it safe to expand a good subarray to $m$'s full window?
  \item What changes between the max version (\probref{c26-goodsub}) and the sum version (LC 907) about ties?
}
\end{recallbox}
